% Defense slides (10 min), IT4033E capstone -- Project Orange (T-ASR).
% Build: cd slides && pdflatex slides.tex && pdflatex slides.tex
% Result tables are the same generated files as the report (../report/tables/), so they stay in sync.
% Numbers written in the text must match report/report.tex (Comments) -- re-check after each refresh.
\documentclass[11pt,aspectratio=169]{beamer}
\usetheme{Madrid}
\usecolortheme{beaver}
\setbeamertemplate{navigation symbols}{}
\usepackage[T5,T1]{fontenc}  % T5: Vietnamese author names
\usepackage{lmodern,amsmath,amssymb,booktabs,tabularx,array,adjustbox}
% Speaker notes (\note{...} in each frame): hidden in slides.pdf, shown by slides-notes.tex
\ifdefined\withnotes
  \usepackage{pgfpages}
  \setbeameroption{show notes on second screen=right}
  \setbeamertemplate{note page}{\vspace{2mm}\small\insertnote}
\fi

\title[T-Adaptive Segment Routing]{Mathematical modelling of T-Adaptive Segment Routing}
\subtitle{Project Orange -- Challenge EURO/ROADEF 2026}
\author[{\fontencoding{T5}\selectfont N.\,T.~Đức, H.\,H.~Chiến}]{{\fontencoding{T5}\selectfont Nguyễn Tiến Đức} (20252076M) \and {\fontencoding{T5}\selectfont Hoàng Huy Chiến} (20261069M)}
\institute{IT4033E Capstone, SoICT}
\date{October 2026}


\begin{document}

\begin{frame}
\titlepage
\note{
Good morning. We are Nguyen Tien Duc and Hoang Huy Chien. Our capstone project is Project Orange, taken from the EURO/ROADEF 2026 challenge: T-Adaptive Segment Routing. In about ten minutes we will describe the problem, our MILP model, how we validated it, and what happened when we solved the twenty instances of set A with an open-source solver.
}
\end{frame}

\begin{frame}{Problem: T-Adaptive Segment Routing}
\begin{itemize}
\item IP backbone $G=(V,A)$, link metric $\omega(a)$, capacity $c(a)$; demands $d=(s_d,t_d)$ with volume $\nu(d,t)$.
\item Packets follow shortest paths; ties are split (ECMP). Congestion = high load $\lambda(a,t)=\text{flow}(a,t)/c(a)$.
\item \textbf{Segment Routing}: each demand gets a list of at most $maxSeg-1$ \emph{waypoints};
shortest-path routing between consecutive waypoints $\Rightarrow$ bend flows away from hot links.
\item Periods $t=0$ (nominal) and $t=1$ (links $q(1)$ down); between periods at most $\kappa(1)$ segment changes.
\end{itemize}
\vspace{2mm}
\begin{block}{Goal}
Choose a segment path per demand and period to minimise link loads.
Experiments: period~0 alone, then periods $\{0,1\}$ with budget, on the 20 instances of set~A.
\end{block}
\note{
The setting is an IP backbone network. Each link has a routing metric and a capacity, and each demand is a source-target pair with a traffic volume. Routers send packets along shortest paths and split ties evenly, which is called ECMP. Congestion is measured by link utilisation: flow divided by capacity.

Segment Routing lets the operator add up to maxSeg minus one waypoints to a demand. Between two consecutive waypoints the traffic still follows shortest paths, so waypoints bend flows away from hot links without changing the metrics.

The T-adaptive part: there are several periods. Period 0 is the nominal network; in period 1 some links are down. Between periods we may change at most kappa segments, so we cannot re-optimise everything from scratch.

Our goal is to choose one segment path per demand and per period that minimises link load. We ran two experiments: period 0 alone, then periods 0 and 1 together with the change budget.
}
\end{frame}

\begin{frame}{Data and split coefficients $r(i,j,a,t)$}
\begin{itemize}
\item Input: \texttt{net.json} ($\omega$, $c$), \texttt{tm.json} ($\nu$), \texttt{scenario.json} ($maxSeg$, $\kappa$, $q$).
\item $r(i,j,a,t)$ = fraction of a unit flow $i\to j$ in $G(t)$ that crosses $a$ (course document):
\[
r(i,j,a,t)=\frac{\sigma_{iu}(t)\,\sigma_{vj}(t)}{\sigma_{ij}(t)}\quad\text{if } d_{iu}+\omega(a)+d_{vj}=d_{ij},\ a=(u,v),\ \text{else } 0
\]
$\sigma$ = number of shortest paths: one Dijkstra with path counting per source, $O(nm\log n)$.
\item Path counting $\neq$ per-node ECMP: toy instance $r(v_0,v_5,a_{1,3},0)=1/3$ vs.\ $1/2$
in the subject's table. We keep the course definition.
\end{itemize}
\note{
The input comes in three JSON files: the network, the traffic matrix, and the scenario with maxSeg, the budget kappa and the failed links.

The key precomputed data is r(i,j,a,t): the fraction of one unit of flow from i to j that crosses link a in period t. Following the course document, it is sigma(i,u) times sigma(v,j) over sigma(i,j), where sigma counts shortest paths, and it is nonzero only if a lies on a shortest path. We compute it with one Dijkstra with path counting per source.

One point we want to be transparent about: path counting is not the same as per-node ECMP splitting. On the toy instance we get one third where the subject's table says one half. We kept the course definition and documented the difference.
}
\end{frame}

\begin{frame}{MILP model}
\textbf{Variables}
\begin{itemize}
\item $x^{dt}_{ij}\in\{0,1\}$: segment $(i,j)\in K$ used by demand $d$ at period $t$ ($K$ = reachable pairs).
\item $U\ge0$: maximum link utilisation (MLU);\quad $y^{dt}_{ij}\ge0$ linearises $|x^{dt}_{ij}-x^{d,t-1}_{ij}|$.
\end{itemize}
\textbf{Objective} $\min U$ \;(first criterion of the subject's lexicographic objective).
\begin{align*}
&\textstyle\sum_j x^{dt}_{ij}-\sum_j x^{dt}_{ji}=\mathbb 1[i=s_d]-\mathbb 1[i=t_d] && \text{(1) path in the segment graph}\\
&\textstyle\sum_{(i,j)} x^{dt}_{ij}\le maxSeg && \text{(2) at most 5 waypoints}\\
&\textstyle\sum_d\sum_{(i,j)} r(i,j,a,t)\,\nu(d,t)\,x^{dt}_{ij}\le U\,c(a)\quad \forall a\notin q(t) && \text{(3) load}\\
&\textstyle\sum_{d,(i,j)} y^{dt}_{ij}\le\kappa(t),\ \ y\ge x^t-x^{t-1},\ y\ge x^{t-1}-x^t && \text{(4) change budget}
\end{align*}
Size: $|D|\,|K|\,|T|$ binaries, $|K|\le n(n-1)$ $\Rightarrow$ from $1.5\cdot10^4$ (inst.\ 01) to $\approx10^9$ (inst.\ 20).
\note{
The variables: x is one if demand d uses segment (i,j) in period t; K is the set of reachable pairs. U is the maximum link utilisation, and y linearises the absolute change between consecutive periods.

We minimise U, the first criterion of the subject's lexicographic objective.

Constraint 1 is flow conservation in the segment graph, so the chosen segments form a path from source to target. Constraint 2 limits the number of segments. Constraint 3 says the load on each working link, built from the r coefficients and the volumes, is at most U times the capacity. Constraint 4 is the change budget between periods.

The weak point is size: binaries = demands times segments times periods. It goes from fifteen thousand on instance 1 to about one billion on instance 20. This drives everything you will see in the results.
}
\end{frame}

\begin{frame}{Implementation and validation}
\begin{itemize}
\item Julia 1.13 + JuMP 1.32 + HiGHS 1.15 (open-source branch-and-bound).
\item \textbf{Validation first}: 24 automated tests on the subject's toy instance; MILP optimum
$=$ brute-force enumeration of all segment paths (MLU $0.375$, waypoint $v_4$).
\item \textbf{MIP start}: no-waypoint routing ($x_{s_dt_d}=1$) is always feasible with 0 changes;
its MLU is the reference column \emph{no-wp}.
\item \textbf{Protocol}: 10 runs per instance (HiGHS seeds 1--10, 1 thread each), time limit 1800~s (period 0) /
3600~s (periods 0+1); gap $=1-\text{LB}/\text{UB}$; CPU = \texttt{optimize!} time only.
\item \textbf{Machine}: i9-14900K, 125~GB RAM (shared server), runs in parallel through a memory-aware job queue.
\end{itemize}
\note{
We implemented the model in Julia with JuMP and the open-source solver HiGHS.

Before any experiment we validated it on the subject's toy instance: 24 automated tests, and the MILP optimum equals brute-force enumeration of all segment paths, MLU 0.375 with waypoint v4.

We give HiGHS a MIP start: routing with no waypoints is always feasible and uses zero changes. Its MLU is the reference column no-wp in the tables.

Protocol: ten runs per instance, seeds 1 to 10, one thread each, time limit 30 minutes for period 0 and one hour for two periods. Gap is one minus lower bound over upper bound; CPU time is the optimize call only.

Everything ran on a shared server, i9-14900K with 125 GB of RAM, through a job queue that only starts a job when its measured memory fits.
}
\end{frame}

\begin{frame}{Results -- period 0 (averages over 10 runs)}
\centering
\adjustbox{max width=\textwidth,max totalheight=0.8\textheight}{\begin{tabular}{lrrrrrrrrrr}
\toprule
Inst & $n$ & $m$ & $|D|$ & vars & opt/runs & MLU & no-wp & gap (\%) & CPU (s) & nodes \\
\midrule
01 & 20 & 80 & 40 & 0.02M & 10/10 & 0.929 & 1.000 & 0.0 & 4.0 & 1.0 \\
02 & 30 & 150 & 45 & 0.04M & 10/10 & 0.549 & 1.000 & 0.0 & 27.8 & 1.0 \\
03 & 50 & 250 & 20 & 0.05M & 10/10 & 0.944 & 1.000 & 0.0 & 16.3 & 1.0 \\
04 & 50 & 250 & 200 & 0.49M & 0/10 & 0.675 & 0.692 & 31.5 & 1819.3 & 0.0 \\
05 & 100 & 396 & 100 & 0.99M & 4/10 & 0.135 & 0.333 & 12.6 & 1578.5 & 0.4 \\
06 & 100 & 500 & 500 & 4.95M & 0/10 & 1.000 & 1.000 & 100.0 & 2056.4 & 0.0 \\
07 & 100 & 500 & 800 & 7.92M & 2/10 & 0.908 & 0.908 & 80.0 & 2069.4 & 0.2 \\
08 & 150 & 654 & 200 & 4.47M & 1/10 & 0.800 & 1.000 & 53.8 & 12450.4 & 0.1 \\
09 & 150 & 750 & 200 & 4.47M & 0/10 & 0.611 & 0.611 & 17.3 & 2403.3 & 0.0 \\
10 & 150 & 966 & 1000 & 22.35M & 0/10 & 1.000 & 1.000 & 100.0 & 2989.6 & 0.0 \\
11 & 200 & 1000 & 400 & 15.92M & 0/10 & 1.000 & 1.000 & 100.0 & 2430.3 & 0.0 \\
12 & 200 & 898 & 400 & 15.92M & 0/10 & 1.000 & 1.000 & 100.0 & 2342.9 & 0.0 \\
13 & 200 & 1000 & 2000 & 79.60M & \multicolumn{6}{c}{model too large} \\
14 & 250 & 1108 & 600 & 37.35M & 0/10 & 1.000 & 1.000 & 100.0 & 3032.5 & 0.0 \\
15 & 250 & 1250 & 600 & 37.35M & 0/10 & 1.000 & 1.000 & 100.0 & 2826.6 & 0.0 \\
16 & 250 & 1452 & 4800 & 298.80M & \multicolumn{6}{c}{model too large} \\
17 & 300 & 1270 & 2000 & 179.40M & \multicolumn{6}{c}{model too large} \\
18 & 300 & 1500 & 2000 & 179.40M & \multicolumn{6}{c}{model too large} \\
19 & 300 & 1998 & 6000 & 538.20M & \multicolumn{6}{c}{model too large} \\
20 & 400 & 2000 & 6000 & 957.60M & \multicolumn{6}{c}{model too large} \\
\bottomrule
\end{tabular}
}
\note{
Period 0, averaged over ten runs. Columns: instance size, number of proved optima, average and best MLU, best bound, gap, time, branch-and-bound nodes, and the no-waypoint MLU for reference.

Read it top to bottom: small instances are solved, medium ones partly, large ones only return the MIP start, and the last ones are marked too large to build. The next slide explains each group.
}
\end{frame}

\begin{frame}{Period 0 -- observations}
% keep in sync with the Comments of report/report.tex
\begin{itemize}
\item \textbf{01--03}: optimal at the root node in seconds; waypoints lower the MLU from 1.00 to 0.93, 0.55, 0.94.
\item \textbf{04--05}: no proof within 30 min, time spent in the root node (0 B\&B nodes).
Strong seed effect: on 05, 4/10 seeds prove the optimum $0.1155$; on 04 only 1/10 improves the MIP start.
$\Rightarrow$ averaging over 10 seeds is necessary.
\item \textbf{06--12} (millions of binaries): the LP relaxation takes most of the time; usually the MIP start is returned.
08: one seed proves $0.319$, 3 stop at $0.561$, 6 keep $1.000$.
10--12, 14, 15: all runs stop with the MIP start and no lower bound (gap 100\%).
\item HiGHS does not always respect the time limit: up to 51940~s for a 1800~s limit (instance 08).
\item \textbf{13--20}: up to $9.6\cdot10^8$ binaries; $\approx2$~GB per million variables $\Rightarrow$ only 14, 15 fit in 125~GB.
\end{itemize}
\note{
Instances 1 to 3 are solved to optimality at the root node in seconds. Waypoints reduce the MLU from 1.00 to 0.93, 0.55 and 0.94; instance 2 shows the strongest gain.

Instances 4 and 5 are not proved within 30 minutes; all the time is spent in the root node. The seed matters a lot: on instance 5, four of ten seeds prove the optimum 0.1155; on instance 4 only one seed improves the MIP start. That is why we average over ten seeds.

From instance 6 on there are millions of binaries. The LP relaxation takes most of the time and usually the MIP start is returned. On instance 8 one seed proves 0.319, three stop at 0.561 and six keep 1.0. Instances 10 to 12, 14 and 15 all stop with the MIP start and no lower bound.

HiGHS does not always respect the time limit: one run on instance 8 took 51,940 seconds for an 1,800-second limit. We report measured times as they are.

Instances 13 and 16 to 20 need up to 960 million binaries. At about 2 GB per million variables, only 14 and 15 fit in 125 GB.
}
\end{frame}

\begin{frame}{Results -- periods 0 and 1 (budget, presolve off)}
\centering
\adjustbox{max width=\textwidth,max totalheight=0.8\textheight}{\begin{tabular}{lrrrrrrrrrr}
\toprule
Inst & $n$ & $m$ & $|D|$ & vars & opt/runs & MLU & no-wp & gap (\%) & CPU (s) & nodes \\
\midrule
01 & 20 & 80 & 40 & 0.05M & 10/10 & 0.929 & 1.000 & 0.0 & 23.4 & 11.0 \\
02 & 30 & 150 & 45 & 0.12M & 7/10 & 0.942 & 1.098 & 5.2 & 2386.5 & 6.1 \\
03 & 50 & 250 & 20 & 0.15M & 3/10 & 0.961 & 1.000 & 24.0 & 3109.1 & 44.4 \\
04 & 50 & 250 & 200 & 1.47M & 0/10 & 0.898 & 1.000 & 32.9 & 3638.6 & 0.0 \\
05 & 100 & 396 & 100 & 2.97M & 0/10 & 0.865 & 1.000 & 77.1 & 5198.7 & 0.0 \\
06 & 100 & 500 & 500 & 14.85M & 0/5 & 1.000 & 1.000 & 100.0 & 3893.4 & 0.0 \\
08 & 150 & 654 & 200 & 13.41M & 0/1 & 1.000 & 1.000 & 68.1 & 3892.2 & 0.0 \\
09 & 150 & 750 & 200 & 13.41M & 0/5 & 1.000 & 1.000 & 66.0 & 3869.9 & 0.0 \\
10 & 150 & 966 & 1000 & 67.05M & \multicolumn{6}{c}{model too large} \\
11 & 200 & 1000 & 400 & 47.76M & \multicolumn{6}{c}{model too large} \\
12 & 200 & 898 & 400 & 47.76M & \multicolumn{6}{c}{model too large} \\
13 & 200 & 1000 & 2000 & 238.80M & \multicolumn{6}{c}{model too large} \\
14 & 250 & 1108 & 600 & 112.05M & \multicolumn{6}{c}{model too large} \\
15 & 250 & 1250 & 600 & 112.05M & \multicolumn{6}{c}{model too large} \\
16 & 250 & 1452 & 4800 & 896.40M & \multicolumn{6}{c}{model too large} \\
17 & 300 & 1270 & 2000 & 538.20M & \multicolumn{6}{c}{model too large} \\
18 & 300 & 1500 & 2000 & 538.20M & \multicolumn{6}{c}{model too large} \\
19 & 300 & 1998 & 6000 & 1614.60M & \multicolumn{6}{c}{model too large} \\
20 & 400 & 2000 & 6000 & 2872.80M & \multicolumn{6}{c}{model too large} \\
\bottomrule
\end{tabular}
}
\note{
The two-period experiment with the change budget, run with presolve off; the next slide explains why. Rows with fewer than ten runs, and rows marked too large, come from time and memory limits, which the report states explicitly.
}
\end{frame}

\begin{frame}{Two periods -- observations}
\begin{itemize}
\item Variables $\times3$ (two copies of $x$ plus $y$).
\item \textbf{Presolve stall}: with HiGHS defaults, instance 04 stayed in presolve $>3$~h for a 1~h limit.
Ablation (60~s limit): removing (4) \emph{or} presolve off $\Rightarrow$ limit respected.
$\Rightarrow$ the dense budget row blocks presolve; all two-period runs use presolve off.
\item Instance 02: link down makes the no-waypoint MLU $1.098$ (congestion); optimum $0.903$ is proved by
7/10 seeds with presolve off vs.\ 2/10 with presolve on.
\item Instance 05: 2/10 seeds find $0.323$ (no-wp $1.000$); the others keep the MIP start.
\item Without presolve memory grows to 3.7--5.5~GB per million variables
$\Rightarrow$ two-period model of 07 ($\approx88$~GB) only alone, 10--20 cannot be built.
\end{itemize}
\note{
With two periods the model has about three times more variables: two copies of x plus y.

First problem: a presolve stall. With default settings, instance 4 stayed in presolve more than three hours with a one-hour limit. A small ablation with a 60-second limit showed that removing the budget constraint, or turning presolve off, each made the limit respected. So the dense budget row blocks presolve, and all two-period runs use presolve off.

Instance 2 is the interesting case: the link failure makes the no-waypoint MLU 1.098, so the network is congested. The optimum 0.903 is proved by seven of ten seeds with presolve off, against two of ten with presolve on.

On instance 5, two seeds find 0.323 against 1.0 without waypoints; the others keep the MIP start.

The price of presolve off is memory: 3.7 to 5.5 GB per million variables. The two-period model of instance 7 needs about 88 GB, so it can only run alone, and instances 10 to 20 cannot be built.
}
\end{frame}

\begin{frame}{Experimental process: trial and error}
\footnotesize
\begin{tabularx}{\textwidth}{@{}>{\raggedright\arraybackslash}p{4.2cm}>{\raggedright\arraybackslash}X>{\raggedright\arraybackslash}X@{}}
\toprule
Symptom & Diagnosis & Fix / consequence \\
\midrule
Toy: $r=1/3$, subject says $1/2$ & path counting $\neq$ per-node ECMP & keep course definition \\
04: no feasible solution in 300~s & stuck in root node & no-waypoint MIP start \\
Laptop 24~GB too small & measured $\approx2$--3~GB / M vars & 125~GB server + memory-aware queue \\
Two-period build 5428~s & \texttt{sum} over JuMP array is quadratic & generator form: 1.4~s \\
\texttt{optimize!} $>3$~h, 1~h limit & ablation: presolve $\times$ budget row & presolve off for two periods \\
Runs far beyond time limit & HiGHS checks limit between steps (likely) & report measured time; reorder queue \\
\bottomrule
\end{tabularx}
\vspace{2mm}
\normalsize
\textbf{Lessons}: validate on a toy instance against brute force first; solver settings are part of the protocol;
for exact MILPs of this size, \emph{memory} is the first limit.
\note{
Since the grading is on scientific method, we logged every problem. The main ones:

The toy split ratio disagreed with the subject: path counting versus ECMP; we kept the course definition.

Instance 4 found no feasible solution in 300 seconds, stuck in the root node: we added the no-waypoint MIP start.

The laptop's 24 GB was too small: we measured memory per million variables, moved to the 125 GB server and wrote a memory-aware queue.

Building the two-period model took 5,428 seconds because sum over a JuMP array is quadratic; the generator form takes 1.4 seconds.

Then the presolve stall, and the runs far beyond the time limit, which we report as measured.

Lessons: validate on a toy instance against brute force first; solver settings are part of the protocol; and for exact MILPs of this size, memory is the first limit, before time.
}
\end{frame}

\begin{frame}{Conclusion}
\begin{itemize}
\item Compact MILP for T-ASR (segment variables, path-counting split ratios), validated against brute force.
\item Exact solving works for small instances and shows large MLU reductions (02: $1.00\to0.55$).
\item Beyond $\approx10^6$ binaries the solver rarely improves the MIP start in the time limit;
beyond $\approx6\cdot10^7$ (period 0) the model cannot even be built in 125~GB.
\item Next steps: restrict candidate waypoints, column generation on segment paths, full lexicographic objective.
\end{itemize}
\vspace{4mm}
\centering\Large Thank you -- questions?
\note{
To conclude: we built a compact MILP for T-adaptive segment routing, with segment variables and path-counting split ratios, validated against brute force.

Exact solving works for small instances and gives large MLU reductions, for example 1.00 to 0.55 on instance 2.

Beyond about one million binaries the solver rarely improves the MIP start within the time limit; beyond about 60 million variables in period 0 the model cannot even be built in 125 GB.

Next steps: restrict the candidate waypoints, column generation on segment paths, and the full lexicographic objective.

Thank you; we are happy to take questions.
}
\end{frame}

\appendix
\begin{frame}{Backup -- periods 0 and 1 with presolve on}
\centering
\adjustbox{max width=\textwidth}{\begin{tabular}{lrrrrrrrrrr}
\toprule
Inst & $n$ & $m$ & $|D|$ & vars & opt/runs & MLU & no-wp & gap (\%) & CPU (s) & nodes \\
\midrule
01 & 20 & 80 & 40 & 0.05M & 10/10 & 0.929 & 1.000 & 0.0 & 32.6 & 1.0 \\
02 & 30 & 150 & 45 & 0.12M & 2/10 & 1.053 & 1.098 & 42.9 & 3224.8 & 0.7 \\
03 & 50 & 250 & 20 & 0.15M & 4/10 & 0.946 & 1.000 & 11.5 & 2840.8 & 0.8 \\
\bottomrule
\end{tabular}
}
\note{
Backup, only for questions: two-period results with HiGHS presolve on, small instances. Compare instance 2: two of ten optimal with presolve on versus seven of ten with presolve off.
}
\end{frame}

\end{document}
